\chapter*{References}
\addcontentsline{toc}{chapter}{References}
\markboth{References}{References}

Every work below is used as a source of ideas, not as material to be copied. The chapter
numbers cited in the \emph{Map to the source books} refer to the chapter numbering of the
edition listed; a reader who wants to check a single claim need only open that edition.

\begin{sloppypar}
\begin{enumerate}[label={[\arabic*]},leftmargin=1.8em]
  \item W. Rudin, \emph{Principles of Mathematical Analysis}, 3rd ed.,
        McGraw-Hill, 1976.
  \item S. Abbott, \emph{Understanding Analysis}, 2nd ed., Springer, 2015.
  \item S. Axler, \emph{Linear Algebra Done Right}, 4th ed., Springer, 2024.
  \item H. Anton, \emph{Elementary Linear Algebra}, 11th ed., Wiley, 2013.
  \item K. Hoffman and R. Kunze, \emph{Linear Algebra}, 2nd ed.,
        Prentice-Hall, 1971.
  \item G. Strang, \emph{Introduction to Linear Algebra}, 6th ed.,
        Wellesley-Cambridge Press, 2023.
  \item J. Stewart, \emph{Calculus: Early Transcendentals}, 8th ed.,
        Cengage, 2016.
  \item J. Stewart, D. K. Clegg, S. Watson, \emph{Calculus: Early
        Transcendentals}, 9th ed., Cengage, 2020.
  \item M. Spivak, \emph{Calculus on Manifolds}, W. A. Benjamin, 1965.
  \item M. Boelkins, D. Austin, S. Schlicker, \emph{Active Calculus:
        Multivariable}, 2nd ed., open text, Grand Valley State University,
        \url{https://activecalculus.org/multi/}.
  \item J. K. Blitzstein and J. Hwang, \emph{Introduction to Probability},
        2nd ed., CRC Press, 2019.
  \item L. Wasserman, \emph{All of Statistics: A Concise Course in Statistical
        Inference}, Springer, 2004.
  \item D. Goldberg, \emph{What Every Computer Scientist Should Know About
        Floating-Point Arithmetic}, ACM Computing Surveys 23(1), 1991.
  \item S. Boyd and L. Vandenberghe, \emph{Convex Optimization}, Cambridge
        University Press, 2004.
  \item D. Bertsimas and J. N. Tsitsiklis, \emph{Introduction to Linear
        Optimization}, Athena Scientific, 1997.
  \item L. A. Wolsey, \emph{Integer Programming}, Wiley, 1998.
  \item I. Goodfellow, Y. Bengio, A. Courville, \emph{Deep Learning},
        MIT Press, 2016.
  \item M. A. Nielsen, \emph{Neural Networks and Deep Learning}, online book,
        Determination Press, 2015.
  \item J. Avigad and P. Massot, \emph{Mathematics in Lean}, version 4.19.0, June
        2025, \url{https://avigad.github.io/mathematics_in_lean/} (also available
        as an Indonesian-language edition).
  \item J. Avigad, L. de Moura, S. Kong, S. Ullrich, \emph{Theorem Proving in
        Lean 4}, Lean FRO, \url{https://lean-lang.org/theorem_proving_in_lean4/}.
  \item I. J. Goodfellow, J. Shlens, C. Szegedy, \emph{Explaining and Harnessing
        Adversarial Examples}, ICLR 2015 (arXiv:1412.6572).
  \item A. Madry, A. Makelov, L. Schmidt, D. Tsipras, A. Vladu,
        \emph{Towards Deep Learning Models Resistant to Adversarial Attacks},
        ICLR 2018 (arXiv:1706.06083).
  \item A. Athalye, N. Carlini, D. Wagner, \emph{Obfuscated Gradients Give a
        False Sense of Security: Circumventing Defenses to Adversarial
        Examples}, ICML 2018 (arXiv:1802.00420).
  \item C. Szegedy, W. Zaremba, I. Sutskever, J. Bruna, D. Erhan, I. Goodfellow,
        R. Fergus, \emph{Intriguing Properties of Neural Networks}, ICLR 2014
        (arXiv:1312.6199).
  \item N. Carlini and D. Wagner, \emph{Towards Evaluating the Robustness of
        Neural Networks}, IEEE Symposium on Security and Privacy 2017
        (arXiv:1608.04644).
  \item J. M. Cohen, E. Rosenfeld, J. Z. Kolter, \emph{Certified Adversarial
        Robustness via Randomized Smoothing}, ICML 2019 (arXiv:1902.02918).
  \item N. Moustafa and J. Slay, \emph{UNSW-NB15: A Comprehensive Data Set for
        Network Intrusion Detection Systems}, Military Communications and
        Information Systems Conference (Milcom), 2015, pp. 1--6.
  \item M. elShehaby and A. Matrawy, \emph{Evasion Adversarial Attacks Remain
        Impractical Against ML-based Network Intrusion Detection Systems,
        Especially Dynamic Ones}, arXiv:2306.05494, 2023.
  \item W. Diffie and M. E. Hellman, \emph{New Directions in Cryptography},
        IEEE Transactions on Information Theory 22(6), 1976, pp. 644--654.
  \item R. L. Rivest, A. Shamir, L. Adleman, \emph{A Method for Obtaining Digital
        Signatures and Public-Key Cryptosystems}, Communications of the ACM 21(2),
        1978, pp. 120--126.
  \item P. W. Shor, \emph{Polynomial-Time Algorithms for Prime Factorization and
        Discrete Logarithms on a Quantum Computer}, SIAM Journal on Computing
        26(5), 1997, pp. 1483--1509.
  \item C. Dwork, \emph{Differential Privacy}, International Colloquium on
        Automata, Languages and Programming (ICALP), LNCS 4052, Springer, 2006,
        pp. 1--12.
  \item M. Abadi, A. Chu, I. Goodfellow, H. B. McMahan, I. Mironov, K. Talwar,
        L. Zhang, \emph{Deep Learning with Differential Privacy}, ACM CCS 2016,
        pp. 308--318 (arXiv:1607.00133).
  \item NIST, \emph{FIPS 203: Module-Lattice-Based Key-Encapsulation Mechanism
        Standard}, National Institute of Standards and Technology, 13 August
        2024, \url{https://csrc.nist.gov/pubs/fips/203/final}.
\end{enumerate}
\end{sloppypar}
